\chapter[My first results chapter]{Long version of my first results chapter title}
    \label{cha:first-results}

\epigraph{%
    The way to get good ideas is to get lots of ideas, and throw the bad ones away.}
    {Linus Carl Pauling}

\begin{remarks}
        Here you can explain meta-information on the origin of this chapter: how was the idea underlying this chapter conceived? Who was involved in the work?
        %
        Has this chapter already published? If so, where (use the cite command!)
        %
        Which parts of the manuscript did you prepare, which parts were written by your collaborators?
        %
        Did you create all of the figures by yourself?
\end{remarks}

% The following space can be used as a "section abstract" where you summarize the work and its main findings (here some parts left as an example)
%
\noindent Example text: Having established the foundations of the Hubbard model and its application as a (self-interaction) correction to \ac{DFT}, we now come to the main outcome of this thesis. 
%
In \cref{cha:figures}, we have seen that Hubbard corrections can markedly improve upon the otherwise poor description by many xc-functionals of strongly localized electrons.
%
Now, we show first results.
%%%%%%%%%%%%%%%%%%%%%%%%%%%%%%%%%%%%%%%%%%%%%%%%%%%%%%%%%%%%%%%%%%%%

\section{Some overview}
    \label{sec:first-overview}

\section{Methodology}
    \label{sec:first-methodology}
    %
    \subsection{Some details}
        \label{subsec:first-details}
    %
    \subsection{Other aspects}
        \label{subsec:first-aspects}

\section{My brilliant idea}
    \label{sec:first-idea}